\documentclass[final,a4paper,12pt]{article}

\usepackage{notes}
\setlength\parskip{1ex}
\setlength\parindent{0cm}
\usepackage{float}
\usepackage{mathtools}
\usepackage{cancel}
\usepackage{amsthm}
\theoremstyle{definition}
\newtheorem{example}{Example}[section]
\newtheorem{problem}{Problem}[section]

\newcommand\qeq{\stackrel{\mathclap{\normalfont\mbox{?}}}{=}}

\title{Geophysics III}
\author{Mixing Tutorial}
\date{}

\begin{document}
\maketitle

\section{Introduction}
Most geologic materials are mixtures of either multiple minerals and/or fluids.
However, when we make geological observations and model the fields we find that
it is not possible to resolve these individual components.  Instead, we are
sensitive to the `average' properties creating or modifying geophysical fields.
Average is a bit of a dangerous word because it can indicate different things in
different situations.  So when someone says average I want you to ask ``What
kind?''.

The typical person refers to the arithmetic mean when they use the term average.
However, there are multiple types of averages that are mathematically defined.
There is a mean, median and mode that you may remember from statistics.  There
are also an aritmetric, geometric, and harmonic means.  Means can be weighted or
unweighted.  This tutorial focuses on each of these averages and when they may
be appropriately applied to Earth's materials.

\section{Statistical average}
There are three averages that are typically used in statistics.  The
mathematical definition may be different for different distribution functions,
but the most commonly used are those for Gaussian distributions.

\section{Pythagorean means}

Aritmetic
\begin{align}
    \text{unweighted:} && \bar{x} &= \frac{1}{n} \sum_{i=1}^{n} x_i \\
    \text{weighted:} &&
        \bar{x} &= \frac{\sum_{i=1}^{n} w_i x_i}{\sum_{i=1}^{n} w_i} \\
\end{align}
Use: computing effective density

Geometric
\begin{align}
    \text{unweighted:} && \bar{x} &= \left( \prod_{i=1}^{n} x_i \right)^{1/n} \\
    \text{weighted:} && \bar{x} &= 
        \left( \prod_{i=1}^{n} x_i^{w_i} \right)^{1/\sum_{i=1}^{n} w_i} \\
\end{align}
Use: computing effective thermal conductivity for random oriented mixtures

The arithmetic and geometric means are mathematically related.  Suppose we have
a set of numbers, $x = (x_1, x_2, \ldots, x_n)$.  We can define a tranform of
$x \mapsto y$ by $y = \log x$.  Arithmetic means
are often appropriate when the values of something are often about the same
order of magintude, but geometric means make a bit more sense when the range of
a property may span orders of magnitude as the do for many types of conductivity
(e.g., electrical and hydraulic conductivity).

Harmonic
\begin{align}
    \text{unweighted:} && \bar{x} &= \left( \sum_{i=1}^{n} x_i^{-1} \right)^{-1} \\
    \text{weighted:} && \bar{x} &= \left( \sum_{i=1}^{n} w_i \right)
        \left( \sum_{i=1}^{n} w_i x_i^{-1} \right)^{-1}
\end{align}
Use: computing effective thermal conductivity for layered conductivity packages

\section{Other common means}
Quadratic mean (RMS)

\section{Additional mixing formulae}
\subsection{Voigt-Reuss-Hill average}
The Voigt-Reuss-Hill average is really an average of the volume weighted
arithmetic and harmonic means and is common for estimating the bulk or effective
magnitude of a physical property.  The average is given by
\begin{align}
    \bar{\Psi} &= \frac{1}{2}(\Psi_V + \Psi_R) \\
    &= \frac{1}{2}\left[
    \sum_{i = 1}^n \Psi_i v_i +
    \left( \sum_{i = 1}^n \frac{v_i}{\Psi_i} \right)^{-1}
    \right],
\end{align}
where $n$ is the number of consituent minerals and $v_i$ is the volume fraction
of the $i$-th component.

\subsection{Hashin-Strikman bounds}
The Hashin-Shtrikman bounds are commonly quoted as a mixing relation for
electrical conductivity and occasionally seismic velocity; however, the bounds
often over- or under-predict the conductivity, but that is what they are
designed to do.  The bounds are
\begin{equation}
    \sigma_{\text{HS}}^- = \mcal{S}(\sigma_{\text{min}}) \leq \sigma \leq
    \mcal{S}(\sigma_{\text{max}}) = \sigma_{\text{HS}}^+,
\end{equation}
where $\mcal{S}$ is defined as
\begin{equation}
    \mcal{S}(x) =
    \Biggl( \sum_{i=1}^N \frac{\phi_i}{\sigma_i + 2x} \Biggr)^{-1} - 2x ,
\end{equation}
and $\phi_i$ is the volume fraction of the $i$-th component.




\end{document}
